\chapter{Fråværet}

\begin{motto}
Eit modent fag kjennest att på kva det har.\\ Eit umode fag kjennest att på kva det manglar, og på at det har slutta å sakne det.
\end{motto}

Vi har no følgt designfaget gjennom heile historia, frå handverket til ekstraksjonsøkonomien, og same mønsteret har dukka opp i kvar epoke: ein lovnad om eit grunnlag, og ein leveranse av noko anna. Dette kapitlet samlar mønsteret i den forma det er lettast å verifisere og vanskelegast å bortforklare. I staden for å spørje kva faget \emph{har}, spør vi kva det \emph{manglar}, og vi gjer det systematisk. For fråværa er ikkje tilfeldige hol. Dei er fire ansikt på den same grunnen som aldri vart lagd, og dei heng saman så tett at kvar av dei føreset dei andre.

\section{Tittelen som kven som helst kan ta}

Byrj med det enklaste å konstatere, fordi det òg har ein juridisk test dei andre fråværa manglar, og fordi testen er knusande. Dei modne profesjonane har organ med makt til å ekskludere, og til sjuande og sist makt til å nekte nokon retten til namnet. Sjå korleis dei gjer det. I Storbritannia fører General Medical Council eit register over legar, og rådet kan stryke ein lege frå det; handlinga heiter på engelsk \emph{erasure}, i daglegtale \emph{striking off}. Er du stroken, er du ikkje lenger lege i lova si meining, og praktiserer du vidare, praktiserer du ulovleg. Advokatstanden har sin parallell i \emph{disbarment}. Og, dette er det mest talande for oss, arkitektane har det same: Architects Registration Board fører eit register, og berre den som står der, har lov å kalle seg «architect» i yrkessamanheng. Tittelen er verna i lov, og rådet kan ta han frå deg \autocite{riba2019code}. Tre felt, tre register, tre måtar å bli kasta ut på. Felles for dei alle er at organet ikkje berre kan klandre deg; det kan slette deg.

No still det same spørsmålet til designen, og høyr stilla. Det finst ingen register over designarar nokon stad. Det finst inga slette-handling, inga \emph{erasure}, inga \emph{disbarment} for ein designar, fordi det ikkje finst noko ein designar kunne strykast frå. Bransjeorgana finst: International Council of Design (tidlegare Icograda, no ico-D), World Design Organization (tidlegare Icsid), AIGA og IDSA i Amerika, og fleire av dei har kodeksar \autocite{icod_code,aiga_standards,idsa_ethics}. Men sjå kva ein slik kodeks faktisk kan gjere. I beste fall kan organet kaste deg ut av ei \emph{frivillig medlemsforeining}. Det kan ikkje frata deg retten til å kalle deg designar, av den enkle grunnen at den retten aldri låg hjå organet. «Designar» er ein ubeskytta tittel i tilnærma kvar jurisdiksjon på jord: kven som helst kan ta han, i dag, utan løyve, utan eksamen, utan register, utan nokon som kan ta han attende.

Det er ikkje eit hol i reguleringa som tida vil tette. Det er ein logisk umoglegheit gjeven faget sin tilstand. Eit ekskluderande organ treng ein standard å ekskludere for brot på. Standarden måtte seie kva ein form skal bere, kva omsyn som er ufråvikelege, kva som skil eit forsvarleg designval frå eit uforsvarleg. Men det er nett ein slik standard faget ikkje har, og kvar gong nokon har prøvd å formulere ein, har faget mint dei på at problema er vrange og kunnskapen situert \autocite{buchanan1992wicked}. Du kan ikkje byggje eit tribunal på «det fungerer ikkje heilt». Fråvêret av organet er ikkje årsaka; det er symptomet. Under det ligg fråvêret av ein standard, og under det igjen fråvêret av eit fundament.

\begin{kasus}{Tre brev og ei foreining}
Set tre personar ved sida av kvarandre. Den fyrste er ein lege som har gjort noko grovt gale; General Medical Council kallar han inn, granskar, og stryk namnet hans frå registeret. Frå den dagen er han ikkje lege. Den andre er ein arkitekt; Architects Registration Board strauk han, og kallar han seg arkitekt på eit prosjekt, bryt han lova. Den tredje er ein designar, kva han enn har gjort: bygd ein mørk mønster-flyt som lurer eldre til abonnement dei ikkje ville ha, ein psykografisk målrettar, ei avhengigheitssløyfe retta mot barn. Kva organ kallar han inn? Kva register står han i? Kva tittel kan takast frå han? Svaret er at ico-D eller AIGA i høgda kan oppheve eit medlemskap han kanskje aldri hadde, og at han morgonen etter framleis er, fullt lovleg, designar. Dei to fyrste breva har adressat og verknad. Det tredje har ingen å sendast til. Det er ikkje at designen er snillare enn medisinen og arkitekturen. Det er at designen ikkje har noko å handheve, og difor ingen til å handheve det.
\end{kasus}

\section{Inga eiga etikk}

Det andre fråvêret følgjer av det fyrste. Faget har inga etikk utleidd av seg sjølv. Det har etikk \emph{lånt}, frå moralfilosofien, frå profesjonsetikken, frå berekraftsdiskursen, og bolta på som eit emne ved sida av det eigentlege faget. Røystene som krev meir, kjem frå randa og taler i imperativ: designarar \emph{bør} ta ansvar \autocite{monteiro2019}, \emph{bør} sentrere rettferd \autocite{costanzachock2020}, \emph{bør} tenkje på framtida \autocite{margolin1998sustainable}. Og appellane er gamle. Manifestet First Things First kom fyrst i 1964, vart skrive om att i 2000, og sa båe gongene at faget brukar kreftene sine på det minst nødvendige \autocite{firstthings2000}. Victor Papanek skreiv i 1971 at design var blitt eit av dei mest skadelege yrka som finst, og bad faget vende seg mot verkelege behov \autocite{papanek1971}. Montreal-erklæringa kom i 2017 og bad om ansvar endå ein gong \autocite{montreal2017}.

At slike appellar må framsetjast, gong på gong, tiår etter tiår, med aukande patos og uendra verknad, alltid utanfrå kjernen, er sjølve diagnosen. I eit fag der etikken fall ut av fundamentet, ville ingen måtte be om han. Økologen treng ingen som ber han hugse at alt heng saman; det er faget hans. Designaren må minnast på det, fordi modellen av faget ikkje inneheld det. Ein etikk som må påleggjast utanfrå, er ikkje fagets eigen; han er ei kledning fagets nakne kropp tek på seg når nokon ser.

Det syner seg klårast på konferansen. Tenk deg ein stor designkonferanse, kva som helst år, kva som helst by: tusen menneske i salen, lanyards og kaffi, ein scene med skarpt lys. Programmet er fullt: ein talar om etisk design og «design for godt», ein om tillit og personvern, ein paneldebatt med tittelen «ansvaret vårt». Folk nikkar, deler sitat på nettet, kjøper boka i pausen. Alt det rette blir sagt. Men det finst ikkje noko tribunal nokon stad i bygget. Det finst ikkje noko rom der ein klage kan leverast, granskast og avgjerast, inga panel med mynde til å seie at den eine talaren, som tilfeldigvis bygde den avhengigheitssløyfa ein annan talar nett fordømte frå scena, ikkje lenger høyrer heime i faget. Begge talar same dag, hentar same applausen. Etikken er på programmet; han er aldri ein instans. Samanlikn med ein medisinsk eller juridisk kongress, der foredraget på scena peikar mot noko verkeleg som kan skje med deg om du bryt: tilsynsorganet, registeret, prosedyren for å bli stroken. På designkonferansen peikar foredraget berre mot fleire foredrag. Difor kan dei rosande orda sirkulere så fritt og kjennast så fine: dei kostar ingenting, fordi dei ikkje er bunde til noko som kan ramme nokon.

\section{Inga delt metode}

Det tredje fråvêret gjeld sjølve grunnlaget for at eit fag kan vere eit fag. Spør om design har ein delt metode, og du finn ikkje berre at svaret er nei, men at faget er usamd om sjølve spørsmålet. Metoderørsla prøvde å gje det ein, og hovudmennene forkasta sitt eige prosjekt. Sidan har det vore god tone å hevde at design \emph{ikkje} har, og ikkje \emph{bør} ha, ein metode, at det er ein «designerly» kunnskapsform \autocite{cross2006}, ein refleksjon-i-handling \autocite{schon1983} som unndreg seg metodisering. Eit fag som ikkje kan einast om hvorvidt det har ein metode, har ikkje ein.

\section{Eit fag utan ordbok}

Og under metoden ligg språket, det fjerde og mest avslørande fråvêret, fordi det er her samanlikninga med dei modne faga blir mest konkret og minst flatterande. Eit fag treng meir enn ord; det treng eit \emph{kontrollert vokabular}, ei mengd termar som tyder det same for alle som brukar dei, slik at ein dom felt av den eine kan etterprøvast av den andre. Medisinen har det i påfallande grad. Han har anatomisk nomenklatur, der kvar muskel, kvart bein, kvar nerve har eitt namn som tyder det same i Oslo og i Tokyo. Han har sjukdomsklassifikasjonar, ICD frå Verdshelseorganisasjonen og den kliniske terminologien SNOMED, som gjev kvar diagnose ein kode, slik at ein journal skriven på den eine sida av jorda kan lesast presist på den andre. Det er ikkje pedanteri. Det er sjølve infrastrukturen som gjer medisinen kumulativ: fordi orda er låste, kan funn leggjast oppå funn utan å skli.

Spør så kva designens kontrollerte vokabular er, og du finn at det ikkje finst. Orda faget feller sine dommar med, «balanse», «integritet», «det fungerer», «rein», «ærleg», er ikkje termar med delte definisjonar. Dei er kalibrert smak. Send ordet på reise og sjå: ein lege i Oslo skriv «hjarteinfarkt» i ein journal, kodar det etter ICD, og sender pasienten vidare til ein lege i Tokyo som aldri har møtt han. Den japanske legen les koden og veit nøyaktig kva som er meint; ordet bar tydinga uskadd over halve kloden, mellom to framande, fordi det er låst til ein definisjon dei begge står under. Send no det andre ordet. Ein designjuryleiar i Oslo skriv at eit arbeid har «god balanse», og sender vurderinga til ein like dugande juryleiar i Tokyo. Kva mottek han? Ikkje ei tyding, men ei stemning. «Balanse» er ikkje låst til noko to framande kunne måle uavhengig og bli samde om; det er eit namn på ein internalisert smak, kalibrert til kvar sine lærarar. Den japanske juryleiaren kan dele kjensla eller ikkje dele henne, og det finst ingen instans over dei to som kan avgjere kven som har rett, fordi det ikkje finst ein definisjon å måle dommane mot. Ordet reiste ikkje; det fordampa ved grensa.

Dette er ikkje fagfolket sin feil. Det er ein direkte følgje av at det aldri vart lagt eit fundament som kunne tvinge orda til å bety det same: det same tomromet Findeli peikte på då han spurde kva designutdanninga eigentleg kviler på \autocite{findeli2001}, og det same Michl avdekte då han synte at faget si grunnsetning var tom \autocite{michl2014}. Og her lukkar ringen seg mot dei andre fråværa. Eit organ som skulle ekskludere, måtte formulere brotet i ord som tydde det same for alle, men orda finst ikkje. Ei etikk som skulle vere fagets eiga, måtte kunne seie presist kva eit forsvarleg val er, men språket til å seie det presist finst ikkje. Ein metode måtte kunne skrivast ned slik at to utøvarar las han likt, men utan eit delt vokabular les dei han kvar sin veg. Språkfråvêret er ikkje eitt fråvær mellom fleire. Det er grunnen dei andre kviler på, eller rettare: grunnen dei alle saknar. Stryk ordboka, og du har stroke føresetnaden for alt det andre på éin gong.

\section{Inga stemme}

Dei fire fråværa kulminerer i eit femte som er summen av dei: faget har inga stemme. Når verda står overfor problem som er formgjevingsproblem i sin kjerne, klimakrisa, den byggde verda, dei digitale systema som omformar barndom og demokrati, er det ikkje designfaget som vert kalla inn som ekspertise, slik klimavitskapen eller folkehelsa vert det. Faget har inga adresse å tale frå, ingen samla posisjon å hevde, ingen kunnskap berre det eig og som verda må gå til det for å få. Det kan stille med einskildpersonar med meiningar. Det kan ikkje stille med eit \emph{fag}.

Og dette, til slutt, er kvifor fråværa er det mest ærlege portrettet av designen. Eit fag som hadde nådd eit fundament, ville hatt alle desse tinga, fordi dei alle følgjer av det same: av å kunne seie presist kva ein form er, kva som verkar på henne, kva ho ber og kva ho legg til side. Det er ein slik presisjon Kuhn meinte skilde det paradigmatiske faget frå det før-paradigmatiske \autocite{kuhn1962}, og det er ein slik presisjon eit moderne fundament for form ville måtte gje \autocite{fields2022}. Designfaget har den ikkje. Og difor manglar det organet, etikken, metoden, språket og stemma, ikkje som fem uavhengige uhell, men som fem målingar av det same fråvêret. Du treng ikkje tru på nokon kur for å sjå det. Du treng berre telje kva som ikkje er der.

\vignett

Men eit fag kan ikkje overleve på fråvær åleine; det treng noko som får fråvêret til å sjå ut som nærvær. Det neste kapitlet handlar om maskineriet som produserer det skinet, akkrediteringa, prisane, fagfellevurderinga, og om kvifor sjølvbedraget er ein stabil likevekt og ikkje ein konspirasjon.

\laeringsmal{\begin{itemize}
\item Gjere greie for skilnaden mellom eit bransjeorgan som kan oppheve eit frivillig medlemskap og eit profesjonsorgan som kan frata ein tittel, og bruke skiljet på medisin, arkitektur og design.
\item Forklare kvifor «designar» er ein ubeskytta tittel, og kva dette syner om fråvêret av ein handhevbar standard.
\item Skilje mellom eit kontrollert vokabular (til dømes anatomisk nomenklatur, ICD, SNOMED) og kalibrert smak, og vise korleis språkfråvêret ligg under dei andre fråværa.
\item Tolke gjentakinga av ansvarsappellane (First Things First 1964/2000, Papanek 1971, Montreal 2017) som evidens for at etikken aldri vart innebygd.
\end{itemize}}

\ovingar{\begin{enumerate}
\item Vel ein lovbeskytta tittel (lege, arkitekt, ingeniør) i din eigen jurisdiksjon. Finn ut kva organ som forvaltar registeret og kva som skal til for å bli stroken. Skriv så ned kva den nærast tilsvarande mekanismen for «designar» er, og forklar kvifor lista blir tom.
\item Ta tre dommar frå ein verkeleg designjury (prisar, opptak, vurdering). Marker kvart vurderande ord, «balanse», «rein», «det fungerer». Prøv å gje kvart ord ein definisjon to framande kunne måle uavhengig. Drøft kva som hindrar deg.
\item Samanlikn ein kodeks frå eit designorgan med éin frå eit profesjonsorgan med handheving (\textcite{icod_code} mot \textcite{riba2019code} eller \textcite{aia_ethics}). Kvar ligg skilnaden, i orda, eller i kva som skjer ved brot?
\item First Things First vart skrive i 1964 og om att i 2000 utan å endre bodskap. Finn ein tredje, nyare ansvarsappell og argumenter for om gjentakinga er framsteg eller symptom.
\end{enumerate}}

\vidarelesing{\fullcite{papanek1971}; \fullcite{findeli2001}; \fullcite{michl2014}; \fullcite{monteiro2019}; \fullcite{costanzachock2020}; \fullcite{margolin1998sustainable}.}
